\properlane{reconciled-multitude}{The Lamb gathers a people who make peace and endure}
\subsection{A numbered people and a multitude beyond counting}
The seer first hears the number of the sealed, then sees the immense multitude. Between those acts of hearing and seeing, the appointment passes over the distribution among the tribes. The remaining contrast is striking: an exact total and a crowd beyond any human count. Bede's explanation of the twelvefold number concerns the whole Church. Its patriarchal descent includes imitation of faith, so membership is wider than natural lineage. His multiplication of the number relates that Church to apostolic preaching throughout the world.\footnote{Revelation 7:4,9; Bede, \work{Explanation}, I, on 7:4, p. 46.}

His account of the following multitude retains a further possibility. After the tribes of Israel, first recipients of the preaching, the text may name the salvation of the Gentiles. This makes the succession historically intelligible as an order of preaching. Bede also gives recapitulation its place: the vision can return to the account of those victorious after persecution. These alternatives concern how the narrative unfolds. They do not compel the conclusion that the multitude outside the numbered group has a lesser share in the Lamb's salvation.\footnote{Bede, \work{Explanation}, I, on 7:9, p. 52.}

The persons in the crowd have nations, peoples and languages. Their common worship does not require the narrative to deny those differences. They stand before the same throne and Lamb and attribute salvation to them. The unity described is consequently a shared relation to the Savior. In this interpretation, the Entrance's common rejoicing with angelic praise is heard beside that worship. Bede explains the sevenfold praise of verse 12 as the Church's confession that its members have received these gifts from God. The congregation's festive praise can be understood in the same dependence.\footnote{Bede, \work{Explanation}, I, on 7:10--12, pp. 52--53.}

The Collect's many intercessors likewise belong to a communion whose members receive from God. Their intercession does not make independent sources of salvation alongside the Lamb. Their shared concern for people still in need becomes explicit again in the Prayer over the Offerings. The difference between the saints' completed life and the living assembly's petition permits assistance across that difference.

\subsection{Why the white robes do not exclude the unbloodied}
The elder identifies the multitude through tribulation and washing in the Lamb's blood. A narrow interpretation might restrict it to those who have shed their own blood in martyrdom. Bede directly addresses that possibility. Martyrs are washed in their own blood, he says, but Jesus' blood cleanses the whole Church from sin. The cause of the multitude's cleansing is thus the saving work of Christ, which extends beyond the manner of any individual's death.\footnote{Bede, \work{Explanation}, I, on 7:14, p. 53.}

That breadth leaves tribulation in place. Bede speaks of entry into the kingdom through many tribulations and emphasizes the final persecution. He gives no date for it. His explanation of standing together concerns faithful confession in adverse circumstances. The robe consequently signifies a received cleansing joined to fidelity. His earlier comment on palms identifies the triumph of the Cross and the victory won in Christ. Neither the garment nor the palm makes suffering its own source of salvation.\footnote{Bede, \work{Explanation}, I, on 7:9,14, pp. 52--53.}

The received Victorinus commentary reinforces the baptismal understanding of the garments and the need to preserve received grace. It differs from Bede at the eastern angel: the received Victorinus text identifies Elias, whereas Bede identifies the incarnate Lord, the messenger of the Father's will. That difference survives their agreement about cleansing. The corporate argument here follows the account of the robes, rather than depending on an agreed identification of every figure.\footnote{Received Victorinus, \work{Commentary}, seventh chapter, comments 2,9; Bede, \work{Explanation}, I, on 7:2, pp. 44--45.}

The psalm begins with the world and everyone who inhabits it belonging to God. Augustine understands the nations believing in the glorified Christ as his Church. The generation of seekers then has a shared moral description: innocent action, a pure heart and desire for God's face. Read beside Revelation, those features explain the life of a gathered people rather than a population assembled by origin alone. Aquinas similarly describes the many who seek God as one generation.\footnote{Augustine, \work{Expositions}, Psalm 23.2--6; Aquinas, \work{Super Psalmos} 23, n. 4, paragraph 86991.}

\subsection{Peacemaking and the Son's work}
Jesus promises that peacemakers will be called God's children. Chrysostom explains the promised name by the work of the only-begotten Son: he unites the divided and reconciles the alienated. Peacemakers resemble him by doing a work of reconciliation. The inference moves from an activity to a family likeness. Sonship names a relationship whose character becomes visible in conduct.\footnote{Chrysostom, \work{Matthew} XV.7.}

He gives the command a wider reach than avoiding one's own quarrels. The disciple must also act among others who are divided. That addition changes the practical meaning of peace. A person can refrain from initiating conflict while remaining indifferent to neighbors whose estrangement continues. Chrysostom requires the further work of helping them become reconciled. The blessing concerns something done, with other persons as its beneficiaries.

Augustine's inward ordering and Chrysostom's reconciliation therefore emphasize different objects. Augustine explains the peace of desires brought under God. Chrysostom requires action between people. Aquinas includes both, with an explicit further step from interior order and peace with others to making peace among those at odds. A person trying to reconcile others still needs his own desires ordered; personal tranquility by itself does not accomplish their reconciliation.\footnote{Augustine, \work{Sermon on the Mount} I.2.9; Aquinas, \work{Super Matthaeum}, cap. 5, lectio 2, paragraph 87340, on 5:9.}

In this interpretation, the Lamb's gathered people and the peacemakers can be considered together: people reconciled in Christ learn to undertake his reconciling work. The inference brings Bede's explanation of the Church's unity into relation with Chrysostom's explanation of the Son's activity. Each commentator supplies a distinct argument. Bede explains who the multitude is and how it is cleansed; Chrysostom explains the conduct that bears the Son's likeness.

The epistle's children already receive their identity from the Father's love. Its wider context joins that identity to love in deeds, including assistance to a brother in need. The work of reconciliation therefore belongs to a life whose source precedes its actions. The saints' intercession, described by the Mass's prayers, gives another example of concern for those still needing help. Their perfected life is compatible with such concern.\footnote{1 John 3:1--3,16--18. Verses 16--18 explain the chapter's context and are not additional appointed verses.}

\subsection{Why peace may bring opposition}
Chrysostom immediately follows his explanation of peacemaking with the next blessing. He says Jesus adds persecution for justice lest every peace be thought blessed. Some agreements would require abandoning virtue or refusing help to others. Chrysostom understands justice broadly here, including virtue, godliness and assistance. Peace has to be tested by those goods. His explanation follows the actual adjacency of the two blessings.\footnote{Chrysostom, \work{Matthew} XV.7.}

The qualification also works in the other direction. Hostility received is insufficient to establish that the recipient acted justly. In the direct address that follows, Jesus specifies false accusations and suffering for his sake. Chrysostom retains both conditions. Aquinas likewise says the cause of persecution makes the difference, and requires falsehood and Christ as the reason in his treatment of verse 11. The moral question concerns what was done and why the person is opposed.\footnote{Chrysostom, \work{Matthew} XV.7--9; Aquinas, \work{Super Matthaeum}, cap. 5, paragraphs 87564--87566.}

This explains why the Communion antiphon's selection matters to the argument. Its three blessings include personal purity, reconciliation and persecution for justice. A social account of holiness has to retain the clean heart. An account of reconciliation has to retain the possibility of faithful opposition. Hearing these promises together allows the communicant to hope for sonship and kingdom without assuming that peacemaking always receives approval.

Chrysostom gives particular attention to reproach because injury often comes without public honor. A person facing danger can sometimes receive praise that encourages endurance; slander can remove even that comfort. His examples explain why Christ's promise addresses the wounds inflicted by false reports as well as bodily harm. The later reference to the prophets belongs to the wider Gospel context; it explains the commentator's reasoning without extending the appointed endpoint.\footnote{Chrysostom, \work{Matthew} XV.8--9.}

John's statement that the world fails to recognize God's children also concerns this difficulty. Augustine explains the failure through Christ's own unrecognized lowliness and his reproof of sin. Recognition by the world therefore cannot decide whether sonship is real. This claim gives no permission to dismiss every criticism. The truth of the charge and the conduct producing it still require the tests Chrysostom derives from Jesus' words.\footnote{Augustine, \work{First John} IV.4.}

\subsection{A community invited to rest and taught to serve}
The acclamation addresses all who labor and are burdened. Chrysostom includes those in anxiety, sorrow and sin. Christ's invitation is extensive in its addressees and personal in its source. Learning his humility concerns the life together of persons who otherwise might compete for honor or cling to their own resentments. The comparison with Revelation is that those gathered receive salvation from Christ; their unity begins in receiving his gift.

The Preface describes the heavenly mother-city and the pilgrims helped by the saints' example and intercession. The people before God and those on pilgrimage can consequently be considered as one communion across different conditions. The mother-city describes belonging and destination. It prevents the multitude's glory from being imagined merely as a crowd of unrelated individual successes. The Prayer after Communion asks concerning perfected charity and the heavenly banquet, so the hope of arrival is also a hope for charity's fulfillment.

If the solemn blessing is selected, its petitions include holy example and service, along with intercession and final joy. Service gives a concrete form to the solidarity developed here. Those who celebrate the saints ask for help to live in a way compatible with their communion. The appointment remains optional; the argument's concern for service already follows from the peacemaking Jesus blesses and the charity for which the Mass prays.

\subsection{The four senses}
\begin{fourSenses}
\item[Literal.] The vision depicts a multitude from the peoples worshipping before God. John names God's children amid the world's failure to recognize them. Jesus blesses those who make peace and those persecuted for justice, specifying false accusation for his sake in the direct address.
\item[Allegorical.] Bede reads the garments through baptism and the palms through Christ's Cross. Chrysostom explains peacemakers through the reconciling Son. Together these accounts explain the identity and activity of a people united to Christ.
\item[Moral.] Christians reconcile divided neighbors, preserve purity and act justly when seeking peace. They test the cause of hostility instead of treating hostility itself as proof of righteousness. The saints' continuing concern for others is compatible with the service asked of those still on pilgrimage.
\item[Anagogical.] The saints stand together before God, and the assembly hopes for a common heavenly banquet in perfected charity. Bede's account includes the whole cleansed Church beyond martyrs. It gives neither an arithmetic population of heaven nor an assurance of each individual's perseverance.
\end{fourSenses}
